% Discussion

This section examines the practical challenges of deploying AnimatorSBT
in the anime industry (\cref{subsec:adoption_barriers}),
discusses the trust model and potential abuse vectors (\cref{subsec:trust}),
considers legal and regulatory implications (\cref{subsec:legal}),
explains why a technological solution is particularly needed
in Japan's anime industry (\cref{subsec:why_japan}),
and explores generalizability to other creative industries (\cref{subsec:generalizability}).

% ----- 7.1 Adoption Barriers -----
\subsection{Adoption Barriers}
\label{subsec:adoption_barriers}

While AnimatorSBT is designed for minimal friction (R5),
several non-technical barriers may impede real-world adoption:

\textbf{Blockchain skepticism.}
The anime industry's workforce skews toward traditional craft skills
rather than digital technology.
The association of blockchain with cryptocurrency speculation and NFT scams
may generate initial resistance.
Our design mitigates this by making the blockchain layer invisible
to end users---animators interact only with the familiar PMC interface.
However, marketing and communication strategies must carefully position
the technology as a ``digital certificate'' rather than an ``NFT'' or ``crypto token''
to avoid triggering negative associations.

\textbf{Studio cooperation.}
While AnimatorSBT can function without studio cooperation
(animators self-report their contributions through the PMC),
studio endorsement would significantly increase the credibility of issued SBTs.
A co-signing mechanism, where studios countersign contribution claims,
would transform the SBT from a self-attestation to a verified credential.
However, studios may resist transparency in subcontracting arrangements,
as it would expose the depth of outsourcing chains
that the industry has historically kept opaque.

\textbf{Network effects.}
The value of AnimatorSBT increases with adoption:
verifiers (studios hiring freelancers) will only check SBTs
if a critical mass of animators holds them.
This creates a chicken-and-egg problem common to credentialing platforms.
We propose addressing this through a phased rollout:
(1)~integrate SBT issuance into the existing PMC user base,
(2)~partner with progressive studios willing to recognize SBTs in hiring,
and (3)~expand to the broader industry as demonstrated value accumulates.

\textbf{Studio incentives.}
AnimatorSBT is designed primarily for animators, but studio cooperation
significantly increases credential value.
Why would studios participate?
We identify three incentive pathways:
(1)~\textit{talent retention}: studios that endorse SBTs signal
transparency and respect for creators, making them more attractive
to high-quality freelancers in a competitive labor market;
(2)~\textit{regulatory compliance}: Japan's Fair Trade Commission
has flagged opaque subcontracting practices~\cite{jftc2025anime},
and SBT co-signing provides auditable evidence of fair attribution;
(3)~\textit{production quality assurance}: verifiable animator portfolios
reduce hiring risk by allowing studios to confirm a freelancer's
actual track record before assignment.

\textbf{Cultural norms.}
Japanese business culture emphasizes relationships and trust
built through personal networks (\textit{kone}).
A cryptographic verification system may be perceived as undermining
this relationship-based hiring model.
We argue that AnimatorSBT complements rather than replaces personal networks:
it provides a baseline of verifiable facts
that enhances (rather than substitutes for) trust-based relationships.

% ----- 7.2 Trust Model and Abuse Prevention -----
\subsection{Trust Model and Abuse Prevention}
\label{subsec:trust}

AnimatorSBT's current trust model relies on self-reported work logs.
This introduces potential abuse vectors that must be addressed:

\textbf{Fabricated contributions.}
An animator could create false work log entries in the PMC
to mint SBTs for work they did not perform.
We propose three mitigation layers:

\begin{enumerate}
  \item \textbf{Studio co-signing (recommended):}
    Extend the smart contract with a two-party confirmation mechanism
    where the studio (or production manager) countersigns the SBT.
    This transforms the SBT from a self-attestation into a mutually verified record.
  \item \textbf{Peer attestation:}
    Allow other animators who worked on the same production to attest
    to a colleague's contribution, creating a social verification layer.
  \item \textbf{Statistical anomaly detection:}
    The Issuance Service can flag statistically improbable claims
    (e.g., an animator claiming 500 cuts in a single week)
    for manual review before minting.
\end{enumerate}

\textbf{Sybil attacks.}
A malicious actor could create multiple wallets to fabricate a diverse portfolio.
Account abstraction mitigates this partially by linking wallet creation
to PMC registration, but a determined attacker could create multiple PMC accounts.
Integration with a decentralized identity protocol (e.g., Proof of Humanity,
Worldcoin) could provide stronger Sybil resistance in future iterations.

\textbf{Issuer key compromise.}
If the Issuance Service's private key is compromised,
an attacker could mint arbitrary SBTs.
The AccessControl pattern mitigates this:
the admin role (held by TOONIQ) can revoke the compromised issuer role
and assign a new key, and all fraudulently minted SBTs can be revoked on-chain.

\textbf{Wallet loss and credential recovery.}
Because SBTs are bound to a specific wallet address,
loss of wallet access means loss of all accumulated credentials.
Account abstraction mitigates this by removing direct private key management,
but the underlying risk remains if the embedded wallet provider
experiences a service disruption.
A social recovery mechanism---where a set of designated guardians
(e.g., trusted colleagues, the PMC operator) can authorize wallet migration
to a new address---is a promising direction.
The smart contract could support a \texttt{migrate()} function
that re-mints all SBTs to a new address upon multi-signature approval,
preserving the credential history.
This is an open problem in the broader SBT literature~\cite{buterin2022decentralized},
with recent work exploring social recovery mechanisms
for SBT-bearing wallets~\cite{goldston2023recovery}.

\textbf{Surveillance and power asymmetry.}
A system that creates persistent, granular records of every task
an animator completes could be repurposed as a surveillance tool.
Studios or clients could use SBT data to monitor productivity
(e.g., ``this animator only completed 20 cuts last month''),
compare animators against each other, or penalize those deemed underperforming.
This risk is particularly acute given the existing power imbalance
between studios and freelance animators~\cite{jftc2025anime},
and mirrors broader concerns about surveillance
in platform-mediated gig work~\cite{sannon2022privacy}.
AnimatorSBT mitigates this through two design choices:
(1)~the animator controls which SBTs are publicly visible
via the selective disclosure mechanism (\cref{subsec:privacy}),
and (2)~metadata granularity can be adjusted---the \texttt{task\_count} field
reports aggregate totals per project, not daily output.
Nevertheless, the tension between verifiability and privacy
is inherent to any credentialing system, and deployment must be accompanied
by clear policies that prevent misuse of contribution data
for punitive performance monitoring.

% ----- 7.3 Legal and Regulatory Considerations -----
\subsection{Legal and Regulatory Considerations}
\label{subsec:legal}

\textbf{Personal data protection.}
AnimatorSBT stores contribution metadata (project title, role, task count, period)
on IPFS with a reference on-chain.
Under Japan's Act on the Protection of Personal Information (APPI),
work history linked to an identifiable individual constitutes personal data.
The use of pseudonymous wallet addresses provides a degree of separation,
but the combination of project title, role, and time period
may be sufficient to identify an individual in practice.
The nullable studio field and selective disclosure mechanisms
described in \cref{subsec:privacy} provide partial mitigation,
but a formal data protection impact assessment is recommended
before production deployment.

\textbf{Token classification.}
AnimatorSBT tokens are non-transferable and carry no monetary value.
Under Japan's Financial Instruments and Exchange Act
and the Payment Services Act,
non-transferable tokens that cannot be traded on secondary markets
are generally not classified as crypto-assets or securities.
However, regulatory clarity on SBTs specifically is still evolving,
and legal counsel should be consulted for production deployment.

\textbf{Labor law implications.}
The Fair Trade Commission's December 2025 report~\cite{jftc2025anime}
found that unilateral price determination without negotiation
may violate the Subcontract Act.
AnimatorSBT could serve as an evidence tool in such disputes,
as the \texttt{evidence\_hash} provides a timestamped record of
work performed and its scope.
However, the legal admissibility of on-chain evidence in Japanese courts
has not been established through precedent.

% ----- 7.4 Why Japan's Anime Industry? -----
\subsection{Why a Technological Solution Is Needed}
\label{subsec:why_japan}

A natural question is why the anime industry requires a blockchain-based solution
when other entertainment industries manage attribution through institutional means.
In Hollywood, labor unions such as the Screen Actors Guild--American Federation
of Television and Radio Artists (SAG-AFTRA),
the Writers Guild of America (WGA),
and the International Alliance of Theatrical Stage Employees (IATSE) enforce credit rights through collective bargaining
agreements. These unions maintain records of members' work history,
negotiate credit arbitration processes, and can call strikes
to enforce compliance---as demonstrated by the 2023 SAG-AFTRA and WGA strikes.

Japan's anime industry lacks comparable institutional infrastructure.
Although JAniCA exists as an industry association,
it functions primarily as a survey and advocacy body
rather than a collective bargaining agent with enforcement power.
The key structural differences are:

\begin{itemize}
  \item \textbf{Unionization rate:}
    The anime industry has no equivalent of SAG-AFTRA's mandatory membership.
    JAniCA membership is voluntary and does not confer negotiating rights.
  \item \textbf{Freelance prevalence:}
    At 47--70\% freelancers~\cite{janica2023survey},
    the anime workforce far exceeds the freelance ratio
    in unionized Western entertainment industries.
  \item \textbf{Credit governance:}
    In Hollywood, unions define credit rules (e.g., the WGA's credit arbitration process).
    In anime, credit decisions rest with studios and production committees
    with no external oversight.
  \item \textbf{Multi-layer subcontracting:}
    While Hollywood outsources VFX,
    anime's gross-subcontracting (\textit{gurosu}) system
    routinely delegates entire episodes through two or three tiers of subcontractors,
    creating deeper attribution gaps.
\end{itemize}

In the absence of institutional solutions,
a technological approach that gives individual animators
verifiable, self-sovereign credentials becomes not merely useful
but structurally necessary.
AnimatorSBT fills the role that unions serve elsewhere:
a persistent, trusted record of who did what.

More broadly, verifiable career records have the potential to elevate
the \textit{social capital} of the animator profession itself.
When animators possess bank-recognizable documentation of their work history,
the downstream effects extend beyond hiring:
mortgage approvals, visa applications, and credit assessments---
all of which require institutional proof of professional activity---become
accessible to a workforce that has historically been unable to produce such evidence.
In this sense, AnimatorSBT is not merely a credentialing tool
but a mechanism for integrating freelance animators
into the broader socioeconomic infrastructure
that salaried workers take for granted.

% ----- 7.5 Generalizability -----
\subsection{Generalizability to Other Creative Industries}
\label{subsec:generalizability}

While AnimatorSBT is designed for the anime industry,
the underlying framework is applicable to any project-based creative industry
where freelancers face similar attribution challenges:

\begin{itemize}
  \item \textbf{Video game development:}
    Game credits face similar issues with outsourced QA testers,
    localization staff, and contract artists often going uncredited.
  \item \textbf{Film and television (live action):}
    Below-the-line crew members, particularly in visual effects,
    frequently work through layers of subcontractors.
  \item \textbf{Music production:}
    Session musicians, ghost producers, and uncredited songwriters
    lack standardized contribution records.
  \item \textbf{Architecture and design:}
    Junior architects and designers contribute substantially
    to projects credited to senior partners.
\end{itemize}

The AnimatorSBT metadata schema (\cref{tab:metadata}) can be adapted
to these domains by modifying the \texttt{role} enum and
adjusting the \texttt{task\_count} semantics.
The smart contract, issuance service, and verification portal
require no domain-specific modifications.

\textbf{International subcontracting.}
Anime production increasingly involves cross-border outsourcing,
with key animation, in-betweening, and coloring tasks
frequently delegated to studios in South Korea, China, Vietnam,
and the Philippines~\cite{mori2011pitfall,tschang2010outsourcing}.
This international subcontracting deepens the attribution gap:
animators working overseas are even less likely to appear in Japanese credits
than domestic subcontractors.
AnimatorSBT is inherently borderless---blockchain credentials
are not jurisdiction-specific, and a wallet address works identically
regardless of the holder's country.
However, cross-border deployment introduces additional challenges:
local data protection regulations (e.g., China's Personal Information Protection Law (PIPL),
the EU's General Data Protection Regulation (GDPR)
for European co-productions), language barriers in metadata,
and the need for local PMC deployment or API integration
with existing production management tools used in each country.
Addressing these challenges is essential for AnimatorSBT
to cover the full scope of modern anime production.
